% !TEX program = pdflatex
% Siemens Industry Software — Software Engineer (EDA / Veloce proFPGA) · Offre 504804 · FR
% Points forts : C++/POO/algorithmes (académique), Linux + Python/Bash, debugging, Git, anglais, Tunis sur site, débutant
% Écart honnête : pas d'expérience prod FPGA/HDL — ne pas inventer ; signaler les transferables + volonté d'apprendre
\documentclass[10pt,a4paper]{article}

\usepackage[french]{babel}
\usepackage[utf8]{inputenc}
\usepackage[T1]{fontenc}
\usepackage{lmodern}
\usepackage[margin=1.05cm]{geometry}
\usepackage{xcolor}
\usepackage{titlesec}
\usepackage{enumitem}
\usepackage{hyperref}
\usepackage{fontawesome5}
\usepackage{microtype}

\setlength{\parindent}{0pt}
\setlength{\parskip}{1pt}
\pagestyle{empty}
\setlist[itemize]{leftmargin=0.85em, itemsep=0.5pt, topsep=0pt, parsep=0pt, partopsep=0pt}

\definecolor{primary}{HTML}{1A365D}
\definecolor{accent}{HTML}{2B6CB0}
\definecolor{muted}{HTML}{4A5568}
\definecolor{rulecolor}{HTML}{E2E8F0}

\newcommand{\TargetTitle}{Ingénieur logiciel \textbar{} C++ · Algorithmes · Debugging · Linux}
\newcommand{\TargetKeywords}{C++ · POO · Algorithmes · Debugging · Linux · Python · Bash · Git · Validation de composants · Intégration système · Anglais}

\hypersetup{
  colorlinks=true,
  linkcolor=accent,
  urlcolor=accent,
  pdfauthor={Hamouda Chekir},
  pdftitle={CV — Hamouda Chekir — Siemens EDA Software Engineer}
}

\titleformat{\section}
  {\color{primary}\normalsize\bfseries}
  {}
  {0pt}
  {}
  [\vspace{0.5pt}{\color{accent}\titlerule[1.4pt]}\vspace{3pt}]
\titlespacing*{\section}{0pt}{7pt}{2pt}

\newcommand{\job}[3]{%
  \textbf{\color{primary}#1}\hfill{\color{muted}\small #2}\\[-1pt]
  {\color{accent}\textit{#3}}\\[1pt]
  \begin{itemize}
}

\newcommand{\proj}[3]{%
  \textbf{\color{primary}#1}\hfill{\color{muted}\small #2}\\[-1pt]
  {\color{accent}\textit{#3}}\\[0.5pt]
  \begin{itemize}
}

\newcommand{\edu}[3]{%
  \textbf{\color{primary}#1}\hfill{\color{muted}\small #2}\\[-1pt]
  {\color{accent}\textit{#3}}\\[2pt]
}

\newcommand{\contactsep}{\hspace{4pt}{\color{rulecolor}|}\hspace{4pt}}

\begin{document}

{\fontsize{20}{24}\selectfont\bfseries\color{primary}HAMOUDA CHEKIR}\\[1pt]
{\normalsize\color{accent}\TargetTitle}\\[3pt]
{\footnotesize\color{muted}%
  \faIcon{envelope}\hspace{2pt}\href{mailto:hamoudachkir@yahoo.fr}{hamoudachkir@yahoo.fr}%
  \contactsep
  \faIcon{phone}\hspace{2pt}+216\,55\,913\,832%
  \contactsep
  \faIcon{map-marker-alt}\hspace{2pt}Tunis, Tunisie (sur site)%
  \contactsep
  \faIcon{linkedin}\hspace{2pt}\href{https://www.linkedin.com/in/hamouda-chekir-5053572b4}{LinkedIn}%
  \contactsep
  \faIcon{github}\hspace{2pt}\href{https://github.com/hamoudachekir}{GitHub}%
}\\[4pt]

\section{PROFIL}
Ingénieur logiciel junior (ESPRIT — TWIN) avec une base solide en \textbf{C/C++, algorithmes et POO}, une pratique quotidienne de \textbf{Linux / Python / Bash}, et un réflexe de debugging rigoureux sur des systèmes multi-composants.
Chez Talan : contribution à la conception, l'intégration et la stabilisation de NextHire (\textbf{16 microservices}), dont des parcours de validation temps réel où la qualité fonctionnelle et l'analyse de cause racine comptaient autant que les nouvelles fonctionnalités.
Anglais courant, à l'aise sous encadrement senior, motivé à progresser vers le logiciel de vérification assistée par le matériel (prototypage FPGA) en montant rapidement sur les outils HDL.

\section{EXPÉRIENCE}

\job{Ingénieur logiciel — Projet de fin d'études}{Fév.\ 2026 -- Août 2026}{Talan \textbar{} Tunis}
  \item Appui aux seniors pour livrer un logiciel réparti sur \textbf{16 microservices} (Node/Express, React, Python FastAPI/Flask) — frontières de services claires, intégration Dockerisée, releases prévisibles.
  \item Renforcé la qualité produit sur un parcours de validation live (WebRTC + contrôles d'intégrité) : traçage des régressions, timing et cas limites jusqu'à un comportement fiable.
  \item Intégré des composants système avec l'équipe logicielle — APIs, contrats de données et modes de panne — pour des livraisons cohérentes sur les plateformes.
\end{itemize}

\job{Stagiaire développeur Full-Stack}{Juin 2025 -- Août 2025}{SmartConseil — Dafe Management \textbar{} Tunis}
  \item Développé et maintenu un back-office opérateur (Laravel, PostgreSQL/MongoDB) ; correction de bugs de production et validation via tests, revues, GitLab CI/CD et Docker.
\end{itemize}

\job{Stagiaire développeur web}{Juin 2024 -- Août 2024}{Big Deal (Club Privilèges) \textbar{} Tunis}
  \item Livré un module de codes cadeaux (Laravel/Filament) de bout en bout et validé le comportement avec Git + Postman avant mise en production.
\end{itemize}

\section{PROJETS SÉLECTIONNÉS}
{\footnotesize
\proj{Cours systèmes \& algorithmes (C/C++)}{ESPRIT}{C · C++ · structures de données · POO · algorithmes}
  \item Formation d'ingénieur en C/C++ : programmation consciente de la mémoire, POO, algorithmes et debugging structuré — base que je souhaite approfondir dans une équipe produit industrielle.
\end{itemize}
\proj{NextHire — validation de composants sous charge}{2026}{Python · parcours temps réel · debugging · Linux · Bash}
  \item Durci un module d'intégrité/validation temps réel pour sessions live ; Python/Bash/Git au quotidien pour reproduire les bugs et automatiser les contrôles — réflexes transférables au scripting EDA (TCL/Bash).
\end{itemize}
}

\section{FORMATION}
\edu{Diplôme d'ingénieur — TWIN (Technologies du Web, de l'Internet et des Réseaux)}{2023 -- 2026}{ESPRIT — École Supérieure Privée d'Ingénierie et de Technologies \textbar{} Tunis}
\edu{Licence en informatique de gestion — E-Business}{2020 -- 2023}{ESEN — École Supérieure d'Économie Numérique \textbar{} Manouba}

\section{COMPÉTENCES}
{\small
\textbf{\color{primary}Cœur de poste :} C, C++ (académique + pratique), algorithmes, POO, debugging, Git, Linux, Python, Bash\\[1pt]
\textbf{\color{primary}Qualité \& intégration :} systèmes multi-composants, traçage de régressions, esprit tests/CI, anglais technique\\[1pt]
\textbf{\color{primary}Aussi utilisé :} Docker, GitLab CI/CD, JavaScript/TypeScript, équipes Agile collaboratives\\[1pt]
\textbf{\color{primary}Montée en compétence (EDA) :} HDL (VHDL/Verilog/SystemVerilog), outils FPGA (Vivado/Quartus), STA — prêt à monter sous mentorat senior
}

\vspace{1pt}
\begin{minipage}[t]{0.70\textwidth}
\section{CERTIFICATIONS}
{\small
\textbf{\color{primary}CCNA — Switching, Routing \& Wireless Essentials} — Cisco (sept.\ 2024)\\[1pt]
\textbf{\color{primary}AWS Academy Graduate — Cloud Foundations} — AWS Academy (mai 2025)
}
\end{minipage}%
\hfill
\begin{minipage}[t]{0.26\textwidth}
\section{LANGUES}
{\small\textbf{Arabe} — Langue maternelle}\\[1pt]
{\small\textbf{Français} — Courant}\\[1pt]
{\small\textbf{Anglais} — Courant}
\end{minipage}

\end{document}
